\begin{multicols}{2}
\begin{description}
    \item[Accesibilidad universal:] condición de no discriminación de cualquier persona, permitiendo la autonomía propia, derechos y disfrute de su entorno, independientemente de tener diversidad funcional o no.
    \item[Ancho libre de paso:] banda de un itinerario peatonal que permite una circulación continua a las personas, con una dimensión mínima de 1,80 metros, sin ningún obstáculo en la totalidad de su recorrido (Orden TMA/851/2021, artículo 5.2.b).
    \item[Deep learning (aprendizaje profundo):] campo del machine learning basado en redes neuronales con múltiples capas, capaz de aprender a reconocer automáticamente elementos complejos en imágenes, optimizando la matriz de pesos del modelo.
    \item[Diseño para todos:] concepto de diseño por el que desde el principio se tiene en cuenta la diversidad de toda la población en materia de utilización y disfrute de los entornos.
    \item[Diversidad funcional:] concepto que define el conjunto de la sociedad con capacidades diversas (física, visual, auditiva, cognitiva o relacional) y que tiene en cuenta a todas las personas independientemente de sus capacidades. Sustituye al término discapacidad, con una connotación más social.
    \item[Extracción geométrica:] procedimiento que utiliza como entrada las máscaras de predicción de la segmentación semántica para obtener medidas reales (como metros) y calcular dimensiones a partir de operaciones vectoriales.
    \item[GitHub:] plataforma basada en Git que actúa como repositorio de código fuente, que se utiliza para el control de versiones y la gestión del ciclo de vida del proyecto.
    \item[GPU (Graphics Processing Unit):] unidad de procesamiento que permite la aceleración por hardware y el cálculo masivo en paralelo, vital para el entrenamiento e inferencia de redes neuronales.
    \item[Ground Truth (verdad de terreno):] conjunto de datos reales de referencia, que pueden obtenerse mediante observaciones de campo, con los que poder evaluar la calidad y precisión de las predicciones de un modelo.
    \item[Itinerario peatonal accesible (IPA):] itinerario con un flujo de paso continuo de al menos 1,80 metros de ancho y 2,20 metros de alto que permite la circulación a cualquier persona de forma segura, cómoda y autónoma (Orden TMA/851/2021, artículo 5.1).
    \item[Kanban:] metodología ágil de gestión de proyectos que utiliza un flujo de tareas por estados (pendiente, en proceso, en pruebas o finalizado), que permite visualizar y organizar el desarrollo de un trabajo o proyecto.
    \item[Movilidad reducida:] persona que tiene una limitación en su movilidad que le afecta en sus desplazamientos de manera funcional. Entre otros, incluye a usuarios/as de silla de ruedas o personas mayores.
    \item[Movilidad reducida temporal:] persona que tiene movilidad reducida durante un periodo transitorio. Se incluyen: personas con muletas, lesiones o familias con carritos de bebé.
    \item[Orden TMA/851/2021:] regulación estatal marco que establece las condiciones mínimas de accesibilidad en vías urbanas, y que define las características mínimas para considerar una vía accesible y la normativa sobre plataformas únicas. Sustituye a la Orden VIV/561/2010.
    \item[Plan de accesibilidad universal (PAU):] instrumento técnico de diagnóstico de la accesibilidad municipal, que debe incluir un informe con el listado de barreras en la vía urbana, priorización de actuaciones y un plan de actuación y presupuestario.
    \item[Producto Mínimo Viable (MVP):] mínima unidad de entrega de un producto que contiene los objetivos fundamentales y esenciales del mismo, que permite ofrecer una funcionalidad mínima en un periodo corto de tiempo.
    \item[Segmentación semántica:] técnica de visión por computador capaz de detectar clases semánticas a partir de una imagen digital (acera, calzada, vehículo, vegetación o mobiliario), clasificando píxel a píxel y utilizando máscaras de predicción.
    \item[Sistema de información geográfica (SIG/GIS):] herramienta tecnológica que gestiona información espacial y cartográfica, permitiendo integrar diferentes capas vectoriales para su análisis y representación en mapas.
    \item[Transfer learning:] técnica del deep learning que reutiliza modelos preentrenados que funcionan bien en tareas generales, que realiza un ajuste de pesos llamado \textit{fine-tuning} para la especialización de un problema específico, optimizando el tiempo de entrenamiento.
    \item[Visión computacional:] subconjunto de la inteligencia artificial especializado en el procesamiento y análisis automático de imágenes digitales, permitiendo la extracción de información semántica estructurada a partir de la naturaleza desestructurada de las imágenes.
    \item[Visor SIG (GIS):] interfaz gráfica interactiva que representa datos en un mapa, permitiendo la visualización y consulta de elementos geográficos al usuario.
\end{description}
\end{multicols}